\section{Part C. Finite residue fields}
\label{sec:C}
We now assume explicitly that $K/\Q_p$ is finite, with residue field
$k=\F_q$, except where C.3 and the implication-table conclusion of C.5
state broader scopes. A superscript $(p')$ on a geometric group leaves the
arithmetic quotient $G_K$ unchanged. These statements extend the questions
on cote 149, manuscript pp.~49--51, scans 99--104, to the finite-residue-field
setting requested in Part C; they are not attributed to those manuscript pages.

\begin{claim}{C.1}{Known}
Under good marked reduction, the outer action on $\Pi^{(p')}$ factors
through $G_k$ and is the outer action on the prime-to-$p$ geometric group of
the special fibre. Arithmetic Frobenius acts on the inertia of a rational
cusp by multiplication by $q$.
\end{claim}
\begin{proof}
SGA 1, XIII, Examples 4.4 and 4.5.3 give the smooth-pair specialization
isomorphism of geometric prime-to-$p$ groups and its functoriality over
the base \citep{BC_SGA1}. Thus the representation is pulled back along
$G_K\to G_k$. At a rational cusp use a local parameter $z$ and the covers
$w^n=z$, with $(n,p)=1$. Inertia acts by $w\mapsto\zeta_nw$;
conjugation by arithmetic Frobenius replaces $\zeta_n$ by $\zeta_n^q$.
Passing to inverse limits gives multiplication by $q$ on
$\prod_{\ell\ne p}\Z_\ell(1)$. With geometric Frobenius the inverse
convention would give $q^{-1}$; the assertion uses arithmetic Frobenius.
\end{proof}

\begin{claim}{C.2}{Known}
For a split stable marked model, wild inertia acts trivially on
$\Pi^{(p')}$ in the outer group, and tame inertia acts by the product of
Dehn twists with node exponents $m_e$. Edge groups and cusp groups inject
as pairwise nonconjugate maximal procyclic subgroups
$\widehat\Z^{(p')}(1)$. Each has a normalizing lift of arithmetic Frobenius
acting by the $q$th power.
\end{claim}
\begin{proof}
Apply the log and semi-graph comparisons of B.3 with
$\Sigma=\{\ell:\ell\ne p\}$. Only Kummer indices invertible in the residue
field occur. Hoshi--Mochizuki's Theorem 4.8 and Lemma 5.4 are stated for
this set of primes, and give the coordinates $m_ec$ of tame inertia
\citep{BC_HMTopics}. The action factors through the tame log base; its
wild kernel is therefore trivial. The node map is
$\{(a,b,c)\in T_{p'}^3:a+b=m_ec\}\rtimes G_k\to T_{p'}\rtimes G_k$.
The injection and commensurable terminality arguments use the same
prime-to-$p$ versions of the cited semi-graph statements. A Frobenius lift
in this node group exists by its surjection to the log base; conjugation
on the antidiagonal is $q$. The cusp calculation is C.1. No assertion of
this kind for the full mixed-characteristic geometric group follows from
this quotient comparison.
\end{proof}

\begin{claim}{C.3}{Proved}
For every hyperbolic curve over the fraction field of a henselian
mixed-characteristic DVR, full-group conditions $b$ and $c$ of scans
101--102 are false. Its full geometric group is not isomorphic to the full
geometric group of any smooth hyperbolic curve over a field of the residue
characteristic, so full-group condition $e$ is false as well. The full
mixed-characteristic table is
\[
\begin{array}{c|ccccc}
 &a&b&c&d&e\\\hline
a&-&\mathrm{no}&\mathrm{no}&\mathrm{yes}&\mathrm{no}\\
b&\mathrm{vac}&-&\mathrm{yes}&\mathrm{vac}&\mathrm{vac}\\
c&\mathrm{vac}&\mathrm{vac}&-&\mathrm{yes}&\mathrm{vac}\\
d&\mathrm{no}&\mathrm{no}&\mathrm{no}&-&\mathrm{no}\\
e&\mathrm{vac}&\mathrm{vac}&\mathrm{vac}&\mathrm{vac}&-
\end{array}
\]
Here ``vac'' means a logically valid implication with an impossible
antecedent in this explicitly stated setting.
\end{claim}
\begin{proof}
If $g\ge1$, $H_1(U_{\bar K},\Q_p)$ surjects onto $V_pJ$ and the latter
has determinant $\chi_p^g$ by the Weil pairing. The cyclotomic character
has infinite image on inertia over such a trait; the corresponding
determinant statement is exactly SGA 7 I, IX, Proposition 5.14(b), whose
hypothesis is a trait with fraction field of characteristic different from
$p$, not a finite extension of $\Q_p$
\citep{BC_SGA7}. Thus $c$ is false. If $g=0$, hyperbolicity gives
$\nu\ge3$. After a finite extension splitting the marks, homology is
$\Q_p(1)^{\nu-1}$, again ramified. Since $b\Rightarrow c$ through
abelianization, $b$ is false too.

For the geometric-group assertion, degree-two cohomology at a prime
$\ell\ne p$ first distinguishes proper from affine groups. In the proper
case the same cohomology and $H^1(-,\F_\ell)$ determine the genus $g$.
The characteristic-zero group's maximal elementary abelian $p$ quotient
has dimension $2g$. In characteristic $p$ the corresponding dimension is
the $p$-rank, at most $g$ (the etale part of the $p$-torsion of its
Jacobian). These cannot agree for $g\ge2$. In the affine case the
characteristic-zero group is topologically finitely generated. The affine
characteristic-$p$ group is not: Artin--Schreier classes give infinitely many
independent elements of $H^1_{\rm et}(V,\F_p)$. To see independence, choose
functions with a single pole of arbitrarily large orders prime to $p$ at
one boundary point; Riemann--Roch provides them for all sufficiently large
orders. Every nonzero $(h^p-h)$ with a pole has leading pole order divisible
by $p$, so a finite nontrivial linear combination with largest chosen pole
order cannot be of that form. Hence no full-group realization exists.

Good marked reduction implies $d$ by C.1, and smooth models give examples
where $a,d$ are true but $b,c,e$ false. The compact-type genus-two example
over $\Q_5$ in B.11 has $d$ true and $a$ false. This proves every
nonvacuous cell. The table is specifically mixed characteristic;
it must not be exported to an unspecified equal-characteristic
trait, where constant families give different full-group behavior.
\end{proof}

\begin{claim}{C.4}{Proved}
Let $V$ be a smooth hyperbolic curve over $\F_q$ and put
$\Sigma=\{\ell:\ell\ne p\}$. A nontrivial closed procyclic subgroup of
$\pi_1(V_{\bar k})^\Sigma$ which is essentially normalized with cyclotomic
action is contained in a cusp inertia subgroup. Consequently the maximal
such subgroups are precisely the geometric cusp groups.
\end{claim}
\begin{proof}
We spell out the adaptation needed in L9 rather than cite Stix's
characteristic-zero formulation as though it had the same hypotheses.
The geometric group and its finite-index subgroups are the pro-$\Sigma$
surface groups of SGA 1, XIII, Corollary 2.12. Their cusp injections and
self-normalizers follow from \citet[Proposition 1.2]{BC_MochComb};
the stronger assertion that distinct cusp groups have trivial intersection
is exactly \citet[Lemma 1.5]{BC_HMNodes}, whose hypotheses permit every
nonempty prime set $\Sigma$. These groups are torsion-free. If an abelian
subgroup $A$ meets a cusp group $I$ nontrivially, then for $a\in A$ the
intersection $I\cap aIa^{-1}$ contains that same nontrivial intersection;
Lemma 1.5 gives $aIa^{-1}=I$, and self-normalization gives $a\in I$.
Thus $A\subset I$. This proves explicitly the prime-to-$p$ form of
Stix's Corollary 13 needed below, including its application to a nonzero
$\Z_\ell$ component of a procyclic subgroup.

Suppose a procyclic subgroup $A$ is not cuspidal and is cyclotomically
normalized. Choose a nonzero $\ell$-component $A_\ell\simeq\Z_\ell$.
The union $W$ of all cusp inertia groups is compact: each is parametrized
by its compact conjugacy orbit and there are finitely many geometric cusp
classes. We have $A_\ell\cap W=1$. The compact set
$A_\ell\setminus\ell A_\ell$ can therefore be separated from $W$ in a
finite quotient of the geometric group. Choose an open normal subgroup
$N$ realizing this separation and stable under an open arithmetic
normalizer; finite generation permits a characteristic refinement.
Set $H=NA_\ell$. The finite cyclic $\ell$-group $H/N$ has the property
that the images of all intersections of cusp groups with $H$ lie in its
subgroup of $\ell$th powers.
Indeed its subgroups are linearly ordered, and an image not so contained
would meet the separated generator coset. Hence the nontrivial degree-$\ell$
quotient $H/NA_\ell^\ell$ kills every cusp inertia of the cover associated
to $H$, and extends over its smooth proper compactification $\bar V'$.

For completeness, this descent is geometric. Choose a normalizer element
$\gamma$ projecting to $\operatorname{Frob}_q^r$, where $r>0$ is such that
the arithmetic normalizer surjects onto $G_{\F_{q^r}}$. The universal
property of $\widehat\Z$ gives a continuous homomorphism
$\widehat\Z\to\overline{\langle\gamma\rangle}$, $n\mapsto\gamma^n$;
its composite with $G_{\F_{q^r}}\simeq\widehat\Z$ is the identity.
The element $\gamma$ normalizes $H$ because it normalizes $A_\ell$ and
$N$ is characteristic. Thus $H\rtimes\overline{\langle\gamma\rangle}$
is an open arithmetic subgroup defining a geometrically connected cover
$V'$ over $\F_{q^r}$. Its smooth compactification is defined over that
same finite field, and conjugation by $\gamma$ induces its Frobenius action
on abelianized geometric homology. The composite
$A_\ell\to H\to T_\ell\Jac(\bar V')$ is nonzero: the unramified
degree-$\ell$ quotient just constructed factors through its reduction
modulo $\ell$ and is nonzero on $A_\ell$. It is equivariant for the
normalizing arithmetic group, acting by $\chi_\ell$ on $A_\ell$.
Thus for some $r>0$ it gives a nonzero map
$\Q_\ell(1)\to V_\ell\Jac(\bar V')$ intertwining Frobenius to the
$r$th power. On the source Frobenius has eigenvalue $q^r$; on the target
all complex absolute values are $q^{r/2}$ by the Weil theorem for the
Jacobian. This is impossible. This is exactly the finite-cover argument
of \citet[Theorem 20, proof]{CF_Stix2008}, with the unavailable characteristic-zero
torus assertion replaced by the explicit finite-field weight calculation.
Every cover used has geometric Galois closure of order supported on
$\Sigma$. Maximality now identifies the full cusp groups.
\end{proof}

\begin{claim}{C.5}{Proved}
For every continuous residue splitting $s:G_k\to G_K$, the prime-to-$p$
outer Frobenius obtained by restriction is realized by a smooth hyperbolic
curve over $k$ if and only if $U$ has potentially good marked reduction.
Allowing a power of Frobenius, or a finite residue extension, gives the same
criterion. This includes extensions needed for good or stable reduction
which are wildly ramified. With $b,c,e$ interpreted using geometric
prime-to-$p$ groups, the implication table is the table of B.8. This last
table of universal implications holds over all henselian traits of residue
characteristic $p$, in equal or mixed characteristic; the realization
criterion preceding it still has the finite-residue-field hypothesis.
\end{claim}
\begin{proof}
If good reduction holds after a finite Galois extension, use the finite
descent datum of B.4 on its good special pair. Prime-to-$p$ specialization
is equivariant for every model isomorphism, including those in wild inertia.
The map to
$\Gal(K_1/K)\times_{\Gal(k_1/k)}G_k$ has kernel $I_{K_1}$, trivial on
the prime-to-$p$ geometric group. Pulling back by $s$ therefore gives
exactly the descended Frobenius action. Thus realization already holds
over $k$.

For the converse take a finite extension with stable marked reduction,
possibly wild, and use the actual restriction square
$E_{U_L}\times_{G_L,s|_H}H=E_s\times_{G_k}H$ from B.4. Enlarge the
residue field to split components, nodes, branches and marks. If the stable
graph has a cycle, the resulting homology has a Frobenius eigenvalue one
from graph homology. A smooth affine or proper curve has only eigenvalues
of absolute values $q^{1/2}$ and $q$ in its geometric homology, so cannot
realize that action. If the graph is a tree, degree-two cohomology first
matches properness. In the proper case a node subgroup contradicts C.4.
In the affine case the multiplicity of the eigenvalue $q$ is $\nu-1$;
the normalized component Jacobians have weight one and contribute no such
eigenvalue. Thus a realizing curve has the same number of cusps. C.2 gives
$\nu+n$ distinct cyclotomic maximal procyclic classes, whereas C.4 gives
only $\nu$. This excludes every nodal stable fibre. The stable fibre is
therefore smooth, proving potential good reduction.

For clarity about wild inertia, once stable reduction holds over $L$ its
wild inertia acts trivially on the prime-to-$p$ group by C.2. Hence the
original wild inertia has finite image, because its intersection with
$G_L$ has finite index. The proof did not presume that image was already
trivial over $K$. Nor did it assume that a given residue splitting restricts
to $G_L$ without the open-subgroup construction.

For the table, Tamagawa gives $a\Leftrightarrow b$, B.7 gives
$c\Leftrightarrow d$, and smooth-pair prime-to-$p$ specialization gives
$a\Rightarrow e$. These implications hold for every henselian trait of
residue characteristic $p$. Counterexamples over local fields suffice to
refute every other universal implication, and can be given at every prime.
The curve $\mathbf P^1_{\Q_p}\setminus\{0,p,1,\infty\}$ has $c,d$ true
and $a,b,e$ false: its rational marked sections give unramified abelian
homology, while its stable model has one separating node of thickness one.
The last assertion follows from the realization criterion just proved.
This construction also works at $p=2$ and $p=3$.

For $e$ true but $a,b,c,d$ false, at an odd prime use
$y^2=x^3-p^2x$ minus infinity. It becomes $Y^2=X^3-X$ over
$\Q_p(\sqrt p)$, a good model because its discriminant is $64$, and
inertia acts by $-1$ on every prime-to-$p$ Tate module. At $p=2$ use
$y^2=x^3+2$ minus infinity. Over $r^2=2$, the substitutions
$x=2X$, $y=r(2Y+1)$ give $Y^2+Y=X^3$, whose discriminant is $-27$.
The nontrivial inertia element acts as $(X,Y)\mapsto(X,-Y-1)$, namely
$[-1]$. This is a concrete example where the good-reduction extension
is wild. The substitutions and discriminants are checked in the attached
program. The compact-type genus-two example in B.11 separately proves
that the abelian failure is not confined to genus-zero boundary collisions.
\end{proof}

\begin{claim}{C.6}{Proved}
Assume split stable marked reduction over $K$, and let $s_x$ be the section
of $x\in U(K)$. Write $H=G_{K(\mu_\infty)}$. Up to the geometric
conjugation determined by the lift of the specialization, the subgroup of
$\Pi^{(p')}$ fixed by $s_x(H)$ is:
\[
\begin{array}{c|c}
\text{specialization of }x& (\Pi^{(p')})^{s_x(H)}\\\hline
\text{smooth unmarked point on }C_v&\Pi_v\\
\text{node }e&\Pi_e\\
\text{marked point }c&I_c.
\end{array}
\]
Here $\Pi_v$ is the injected prime-to-$p$ fundamental group of the
normalized component punctured at all its branches and marks. At a node
$xy=\varpi^m$ the contact integers are $a=v(x)$, $m-a=v(y)$ with
$0<a<m$; at a mark the contact integer is $\mu>0$. The listed fixed
groups do not change when these positive integers change.
\end{claim}
\begin{proof}
All prime-to-$p$ roots of unity generate the maximal unramified extension
of a local field: the residue fields $\F_{q^r}$ are generated by roots of
unity of order $q^r-1$. Therefore $H\subset I_K$. Its image in tame inertia
$T_{p'}$ is an open nonzero subgroup. Indeed the remaining extension by
$p$-power roots of unity has a pro-$p$ Galois subgroup of finite index,
so its prime-to-$p$ tame quotient is finite. Wild inertia acts trivially
also on the genuine pointed action: an outer-trivial action gives, by
centre-freeness, a homomorphism from the pro-$p$ group to the pro-prime-to-$p$
group of inner automorphisms, which must be trivial.

Work in the log extension $\Pi_I=\Pi^{(p')}\rtimes^{\rm out}T_{p'}$.
For a point reducing smoothly away from marks, the image of inertia under
its section is the vertex inertia $I_v$, the subgroup centralizing the
vertex group. This follows by the local smooth chart: its generic covers
restrict to the ordinary covers of the punctured normalized component.
The image $J$ of $H$ is a nontrivial closed subgroup of $I_v$.
Hoshi--Mochizuki's Lemma 3.6(i) states, for every such $J$, the exact
equality
$Z_{\Pi_I}(J)\cap\Pi^{(p')}=\Pi_v$
\citep[Lemma 3.6]{BC_HMNodes}. Its hypotheses are the nodally nondegenerate
semi-graph representation; positive thicknesses give the stricter SNN
condition here by their Lemma 2.11 and the length calculation. This proves
the first row, including the smooth-fibre case $\Pi_v=\Pi^{(p')}$.

At a node the local inertia image has coordinates
$(a\tau,(m-a)\tau,\tau)$ in the rank-two group of B.3, and units do not
alter tame inertia. It lies in the edge decomposition group $D_e$ but in
neither adjacent vertex inertia: those correspond respectively to one
branch coordinate zero. Multiplication by a positive integer is injective
on $T_{p'}$, so the same remains true after restriction to $J$.
If $h\in\Pi^{(p')}$ centralizes $J$, then
$J\subset D_e\cap D_{he}$. Proposition 3.8(i)--(ii) of
\citet{BC_HMNodes} says that distinct edge decomposition groups have zero
intersection unless the edges share a vertex, in which case their
intersection is exactly that vertex inertia. Both alternatives contradict
the preceding observation. Hence $he=e$ and $h\in\Pi_e$ by edge
self-normalization. The rank-two group is abelian, so every element of
$\Pi_e$ does centralize $J$. This proves equality.

At a marked point the same proof uses its local group
$I_c\times T_{p'}$ and slope $(\mu\tau,\tau)$. The adjacent vertex inertia
has first coordinate zero; the positive slope excludes it. Proposition
3.8 applies to cusp edges as well as nodes, giving exactly $I_c$.
All subgroup comparisons are made in the common log extension, so the
answer concerns fixed elements of the genuine pointed action, not just
fixed conjugacy classes.
\end{proof}

\begin{claim}{C.7}{Proved}
For $U=\mathbf P^1_{\Q_5}\setminus\{0,1,5,\infty\}$ and $x=2$, choose
geometric cusp generators $a_0,a_5,a_1$ in the free pro-prime-to-five group.
Then
\[
 (\Pi^{(5')})^{s_2(G_{\Q_5(\mu_\infty)})}
 =\overline{\langle a_0a_5,a_1\rangle},
\]
a free pro-prime-to-five subgroup of rank two. For good reduction the
corresponding fixed prime-to-$p$ group at an integral unmarked point is
all of $\Pi^{(p')}$.
\end{claim}
\begin{proof}
Blow up the collision of $0$ and $5$ on the marked projective line over
$\Z_5$. The two resulting rational components each have three special
points and meet in one node of thickness one. The point $2$ lies on the
component with $1,\infty$, and the node. Its node boundary word is
$a_0a_5$, so the normalized component group is the displayed free group.
C.6 proves the equality. The genuine tame action can be written
$a_i\mapsto d a_i d^{-1}$ for $i=0,5$, $a_1\mapsto a_1$, with
$d=a_0a_5$; the code verifies the fixed boundary words and a finite quotient
on which this action is nontrivial. The good-reduction assertion is the
node-free case of C.6.
\end{proof}

\begin{claim}{C.7u}{Unresolved}
The full-group fixed subgroups requested in C4, including the point $2$
and the integral tangent at $0$ on
$\mathbf P^1_{\Q_5}\setminus\{0,1,\infty\}$, are not determined here.
\end{claim}
The quotient map
$\Pi\twoheadrightarrow\Pi^{(p')}$ sends fixed elements to fixed elements,
but invariants need not surject through a quotient. In particular the large
fixed quotient just computed supplies neither a lift nor an upper bound
for the full fixed group. Wild geometric covering data and its genuine
action remain to be controlled; no claim of full-group triviality or cusp
equality is made from the prime-to-$p$ calculation.

\begin{claim}{C.8}{Proved}
After choosing a place of $\bar\Q$ above $p$ and identifying its decomposition
group with $G_{\Q_p}$, its intersection with $G_\Q^\circ$ is
$G_{\Q_p(\mu_\infty)}$. Thus C.6 describes the local cyclotomic-kernel fixed
group at that chosen place in the stated quotient.
\end{claim}
\begin{proof}
The chosen place embeds the algebraic closure into an algebraic closure of
its completion. Finite extensions of the henselization and completion
have the same Galois category, identifying the decomposition group with
$G_{\Q_p}$. Cyclotomic characters commute with this embedding because
their actions are literally the actions on the same roots of unity.
Taking kernels therefore gives the asserted intersection. Another place
changes the decomposition subgroup by $G_\Q$-conjugacy, the stated ambiguity.
\end{proof}

